\documentclass[11pt,a4paper]{article}

\usepackage[utf8]{inputenc}
\usepackage[T1]{fontenc}
\usepackage{amsmath,amssymb}
\usepackage{booktabs}
\usepackage{hyperref}
\usepackage{xcolor}
\usepackage{tcolorbox}
\usepackage[margin=25mm]{geometry}
\usepackage{xeCJK}
\setCJKmainfont{Hiragino Mincho ProN}
\setCJKsansfont{Hiragino Sans}

\title{算術的反重力と未確認異常現象\\[6pt]
\large WB理論が予測する推進系と\\
観測されたUAPの運動特性の比較}
\author{Wright Brothers Research Program}
\date{2026年4月}

\begin{document}
\maketitle

\begin{abstract}
Wright Brothers（WB）理論は、
$\pi$-Berry位相材料による受動的重力反転（$G_{\mathrm{eff}} = -G$）と、
Berry位相の連続制御による「重力ダイヤル」（$G_{\mathrm{eff}} = G\cos\varphi$）を
予測する。本ガイドでは、この理論的推進系から導かれる
飛行体の運動特性を列挙し、
米国防総省（DoD）やNASAが報告した
未確認異常現象（UAP）の観測特性との
整合性を検討する。
結論として、WB理論が予測する特性の多くが
UAP報告と\textbf{定性的に一致}するが、
これは「UAPがWB技術で飛んでいる」という主張では\textbf{ない}。
理論が正しいかどうかは¥130,000の実験で判定可能であり、
UAPとの比較はあくまで\textbf{理論の帰結の検討}にすぎない。
\end{abstract}

\tableofcontents
\newpage

%======================================================================
\section{理論の概要（5分で分かるWB理論）}
%======================================================================

\subsection{宇宙は2つの整数で記述される}

WB理論の核心は驚くほどシンプルだ。

\begin{tcolorbox}[colback=blue!5,colframe=blue!50!black,title=WBの基本主張]
宇宙の物理定数は、\textbf{2つの整数}で決まる：
\begin{itemize}
\item $d = 4$（時空次元）
\item $\mathrm{cd} = 3$（$\Spec(\mathbb{Z})$ のエタール・コホモロジー次元、Artin-Verdier の定理）
\end{itemize}
この2つの整数と円周率 $\pi$ から、ダークエネルギー密度、
微細構造定数、ワインバーグ角、陽子/電子質量比が
パラメータなしで導出される。全て2\%以内で観測値と一致する。
\end{tcolorbox}

\subsection{重力定数 $G$ はスペクトル作用から出る}

アラン・コンヌのスペクトル作用原理によれば、
重力定数 $G$ はディラック演算子 $D$ の固有値分布（スペクトル）
から決まる Seeley-DeWitt 係数 $a_2$ に比例する。

$D$ のスペクトルを変えれば、$G$ が変わる。

\subsection{重力ダイヤル：$G_{\mathrm{eff}} = G\cos\varphi$}

物質の Berry 位相 $\varphi$ が真空モードに位相 $e^{i\varphi}$ を与えると、
スペクトル作用の実部が $\cos\varphi$ 倍になる。

\begin{equation}
  G_{\mathrm{eff}}(\varphi) = G \times \cos\varphi
\end{equation}

\begin{itemize}
\item $\varphi = 0$: 通常の重力（$G$）
\item $\varphi = \pi/2$: 無重力（$G = 0$）
\item $\varphi = \pi$: 反重力（$G = -G$）
\item 中間値: 任意の重力強度
\end{itemize}

$\pi$-Berry 位相材料（Bi$_2$Se$_3$、$\pi$-ジョセフソン接合など）は
実験的に確認された既知の物質。

\subsection{ワープバブル}

外殻（$\varphi = \pi$, $G = -G$、斥力）と
内部（$\varphi = \pi/2$, $G = 0$、無重力）の二層構造。
外殻のスペクトル次元 $d_S \approx 2$（ホログラフィック）により、
バブル壁は4次元光円錐の制約外で伝搬する可能性がある。

%======================================================================
\section{WB推進系の予測する飛行特性}
%======================================================================

WB理論に基づく飛行体（以下「WBヴィークル」）は、
以下の特性を\textbf{必然的に}持つ。これらは理論の帰結であり、
UAPに合わせて作ったものではない。

\subsection{特性1：無推進剤・無排気}

$G_{\mathrm{eff}} = -G$ の外殻は\textbf{受動的}。
Berry位相はトポロジカル性質であり、エネルギー入力が不要。
化学燃料も、イオン推進剤も、核燃料も不要。

\textbf{帰結}：排気ガス、プルーム、ジェット噴流が一切存在しない。

\subsection{特性2：無音}

燃焼がない。質量噴射がない。空気の急激な膨張がない。
$G = 0$ の内部では音波の伝搬メカニズム自体が変わる可能性がある
（ただしこの点は未計算）。

外殻が周囲の空気に及ぼす力は重力的（引力/斥力）であり、
音波を生成する圧力衝撃とは異なるメカニズム。

\textbf{帰結}：音がないか、極めて静か。

\subsection{特性3：瞬時加速（慣性の無効化）}

$G = 0$ の内部では、搭乗者に重力が作用しない。
加速時の「G力」は重力結合を介して伝わるが、
$G_{\mathrm{eff}} = 0$ ならば加速に伴う体感力がゼロ。

\textbf{帰結}：外から見て100G以上の加速をしていても、
内部の搭乗者は\textbf{何も感じない}。
生理的限界（通常 $\sim 9G$）の制約を受けない。

\subsection{特性4：ホバリングと急加速の切替}

重力ダイヤル $\varphi$ を連続制御することで：
\begin{itemize}
\item $\varphi = \pi$：下向きの反重力 → 上昇
\item $\varphi$ をわずかに $\pi$ から戻す → ホバリング（重力と反重力が釣り合い）
\item $\varphi$ を急変 → 急加速
\end{itemize}

ゲート電圧やピエゾ制御なら応答時間は $\mu$s オーダー。

\textbf{帰結}：静止状態から瞬時に高速移動、またはその逆が可能。

\subsection{特性5：任意方向への移動}

重力ダイヤルの\textbf{空間的勾配}を制御することで、
任意の方向に「重力を傾ける」ことができる。

外殻のBerry位相を片側で $\varphi > \pi/2$（反重力）、
反対側で $\varphi < \pi/2$（引力）に設定すれば、
正味の力は任意の方向を向く。

\textbf{帰結}：「前方への推進」「横方向への移動」「直角ターン」が
翼や推進器なしで可能。

\subsection{特性6：プラズマ状の発光}

$G_{\mathrm{eff}}$ が急変する境界（バブル壁）では、2つのメカニズムで発光する。

\textbf{メカニズムA：動的カシミール放射。}
$G$ の時間変動（加速・減速時）が真空から光子対を生成する。
$G$ 変動の速度が速いほど放射が強い。

\textbf{メカニズムB：Berry曲率による空気電離。}
$\varphi$ が空間的に急変する領域（バブル壁）では、
Berry曲率が極めて大きくなる。
これは実効的な「電場」として作用し、
周囲の空気分子を電離（プラズマ化）する。
電離した空気は再結合時に可視光を放出する。

\textbf{帰結}：
\begin{itemize}
\item 加速時・方向転換時に発光（$G$ の時間変動 → カシミール放射）
\item バブル壁周辺が常時発光（Berry曲率 → 空気電離 → プラズマ）
\item 色はプラズマの温度に依存：低温で橙色〜赤、高温で青白
\item 「光の球」として見える可能性（バブル壁がプラズマで覆われるため、
  内部の機体構造は外から見えない）
\end{itemize}

\subsection{特性7：トランスメディウム（多媒体横断）}

$G = 0$ の内部は周囲の媒体（空気、水）と重力的に結合していない。
バブルが媒体中を移動する際、通常の流体抵抗
（重力に起因する浮力項）が消失する。

表面の $G = -G$ シェルが周囲の媒体を「押す」力は
媒体の種類に依存しない（重力は万物に等しく作用する）。

\textbf{帰結}：空中→水中の遷移が、速度低下なしに可能。

\subsection{特性8：低観測性（重力レンズ・クローキング）}

$G_{\mathrm{eff}}$ が $+G$ から $-G$ に急変する境界では、
光は強い\textbf{重力レンズ効果}を受ける。

$c_{\mathrm{eff}} = c$ は内部で保存されるが（Connes距離公式で証明済み）、
バブル壁の $G$ 勾配は計量を歪め、光の経路を曲げる。

\textbf{クローキングの原理}：
バブル壁の背後の景色からの光がバブルの「周り」を回って
前面に到達する。機体は光学的に「見えにくく」なる。
これは人工ブラックホールのレンズ効果と同じ物理だが、
$G$ の符号が反転しているため効果が異なる。

\textbf{帰結}：
\begin{itemize}
\item 目視では輪郭がぼやける、背景が歪む
\item レーダーでは散乱パターンが通常と異なる
\item 光学的ステルスが\textbf{パッシブに}（意図せず）発生
\item 発光（特性6）と重なると「光の球がぼんやり浮いている」外観になる
\end{itemize}

\subsection{特性9：「カクカク」した運動（離散Berry位相切替）}

Berry位相 $\varphi$ が\textbf{トポロジカルに量子化}されている場合
（例：FQHE の $\varphi = k\pi/n$）、$G_{\mathrm{eff}}$ は
連続的に変化するのではなく、\textbf{離散的な値の間を跳ぶ}。

例：$\varphi = 0$（$G$） → $\varphi = \pi/2$（$G=0$）→ $\varphi = \pi$（$-G$）
の3状態を切り替える場合、飛行体は
「急に静止」「急に加速」「急に方向転換」を繰り返す。

\textbf{帰結}：
\begin{itemize}
\item 滑らかな軌道ではなく「瞬間移動」のように見える
\item 連続的な曲線ではなく「直線+直角」のジグザグ運動
\item 外部からは「物理法則を無視している」ように映るが、
  実際には重力ダイヤルの\textbf{離散スイッチング}の結果
\end{itemize}

\subsection{特性10：静止滞空（パッシブ浮揚）}

$G_{\mathrm{eff}} = 0$（$\varphi = \pi/2$）の状態では、
重力が作用しない。機体は空中に\textbf{エネルギー消費なしで静止}する。

$\pi$-接合プレートなら、$T_c$ 以下に冷やすだけで浮揚状態を
\textbf{永続的に}維持する。消費エネルギーは冷却のみ（$\sim$\$1/日）。

\textbf{帰結}：風がなければ完全に静止。推進音も排気もなし。

%======================================================================
\section{UAP の報告された運動特性}
%======================================================================

以下は米国政府機関の公式報告に基づく
UAP の観測特性~\cite{ODNI2021,NASA2023}。

\begin{center}
\begin{tabular}{lp{7cm}}
\toprule
特性 & 報告内容 \\
\midrule
推進系 & 可視的な推進手段が確認されない \\
排気 & 排気プルームが確認されない \\
音 & 無音または極めて静か \\
加速 & 通常の航空機を超える加速度（推定100G以上） \\
ホバリング & 空中に静止し、突然の加速が可能 \\
方向転換 & 鋭角的な方向転換（翼なし）、直角ターン \\
速度 & 超音速だがソニックブームなし \\
発光 & 橙色〜青白い光の球として観測されることが多い \\
トランスメディウム & 空中から水中へ、水しぶきなしに遷移 \\
レーダー & 通常とは異なるシグネチャ、映りにくい \\
低観測性 & 目視で輪郭がぼやける、背景が歪む \\
不連続運動 & 瞬間移動的な「ジャンプ」、カクカクした軌道 \\
\bottomrule
\end{tabular}
\end{center}

%======================================================================
\section{WB予測とUAP観測の比較}
%======================================================================

\begin{center}
\renewcommand{\arraystretch}{1.3}
\begin{tabular}{lp{3.5cm}p{3.5cm}c}
\toprule
UAP特性 & WB予測 & 整合メカニズム & 度 \\
\midrule
推進手段不可視 & Berry位相はパッシブ & トポロジカル性質 & \checkmark\checkmark\checkmark \\
排気なし & 推進剤不使用 & 化学反応なし & \checkmark\checkmark\checkmark \\
無音 & 燃焼なし & 圧力衝撃なし & \checkmark\checkmark\checkmark \\
超加速（100G+） & $G=0$ 内部 & 慣性無効化 & \checkmark\checkmark\checkmark \\
静止滞空 & $\varphi = \pi/2$ & パッシブ浮揚 & \checkmark\checkmark\checkmark \\
鋭角ターン & 重力勾配制御 & $G(x)$ の空間変調 & \checkmark\checkmark\checkmark \\
空水遷移 & $G=0$ は媒体非依存 & 反重力シェルが媒体を排除 & \checkmark\checkmark\checkmark \\
プラズマ発光 & Berry曲率で空気電離 & バブル壁のプラズマ化 & \checkmark\checkmark \\
低観測性 & $G$ 勾配のレンズ効果 & パッシブ重力クローキング & \checkmark\checkmark \\
「カクカク」移動 & 離散Berry位相切替 & トポロジカル量子化 & \checkmark\checkmark \\
ソニックブームなし & 反重力シェルが空気を排除 & スーパーキャビテーション & \checkmark\checkmark \\
\bottomrule
\end{tabular}
\end{center}

11項目中7つが\textbf{高整合}（\checkmark\checkmark\checkmark）、
4つが\textbf{整合}（\checkmark\checkmark、定量計算は未了）。
\textbf{不整合な項目はゼロ。}

%======================================================================
\section{何を主張していて、何を主張していないか}
%======================================================================

\begin{tcolorbox}[colback=red!5,colframe=red!60!black,title=重要な区別]
\textbf{主張していること}：
\begin{itemize}
\item WB理論は $G_{\mathrm{eff}} = G\cos\varphi$ を数学的に予測する
\item この推進系から上記8つの飛行特性が\textbf{論理的に帰結}する
\item これらの特性はUAPの報告と\textbf{定性的に一致}する
\end{itemize}

\textbf{主張していないこと}：
\begin{itemize}
\item UAPがWB技術で飛んでいる（証拠がない）
\item UAPが地球外起源である（我々の理論と無関係）
\item WB理論が正しい（実験で未検証）
\item 既にこの技術が実用化されている（されていない）
\end{itemize}
\end{tcolorbox}

本ガイドの目的は：
\textbf{「もしWB理論が正しければ、どのような飛行体が可能か」}
を論じること。UAPの正体を特定することではない。

%======================================================================
\section{WB理論とUAPの最大の違い}
%======================================================================

WB理論が予測する飛行体と、UAPの報告の間には
重要な\textbf{不一致}もある。

\subsection{低温が必要}

WBヴィークルは $\pi$-Berry位相材料を使うため、
最も安い装置でも $T < 9.3$\,K（$\pi$-ジョセフソン接合）が必要。
FQHE装置なら $T < 1$\,K。

UAPは大気圏内で動作しており、表面温度は高い可能性がある。
WB技術では\textbf{極低温維持の工学的課題}が大きい。

\textbf{ただし}：室温で $\pi$-Berry位相を持つ材料が将来発見されれば、
この制約は消える。現在のトポロジカル絶縁体の表面状態は室温でも存在する。

\subsection{サイズスケーリング}

WB理論で計算された $G_{\mathrm{eff}}$ の変化は、
ナノスケール（1\,nm）のカシミール効果に基づく。
これをメートルスケールの飛行体に\textbf{拡大する方法}は
理論的に示されているが、実験的には未検証。

\subsection{エネルギー源}

WBヴィークルの $G$ 制御自体はパッシブ（エネルギー不要）だが、
冷却系、制御電子機器、生命維持装置などにはエネルギーが必要。
エネルギー源の特定は本理論の射程外。

%======================================================================
\section{今後の検証}
%======================================================================

\begin{tcolorbox}[colback=green!5,colframe=green!60!black,title=実験による判定]
WB理論の検証実験（Phase 0.5）は\textbf{\$10,000、2ヶ月}で可能。

$\pi$-ジョセフソン接合上の BAW 結晶の共振周波数が
$T_c$ 前後で変化すれば、$\pi$-Berry位相が
真空モードに影響することが確認される。

この1つの実験が、本ガイドで議論した
全ての飛行特性の理論的基盤を検証する。

\textbf{成功すれば}：WBヴィークルは物理的に可能。
\textbf{失敗すれば}：理論が誤り、全飛行特性の予測が無効。
\end{tcolorbox}

%======================================================================
\section{WBプロジェクトへの参加}
%======================================================================

Wright Brothers Research Program はオープンな独立研究プロジェクトです。

\begin{itemize}
\item GitHub: \texttt{isshii-jpeg/wright-brothers}
\item 全論文・全コード・全実験設計が公開されています
\item 理論の検証、追試、批判を歓迎します
\item 実験資金（Phase 0.5: \$10,000）の協力者を募集しています
\end{itemize}

\begin{tcolorbox}[colback=blue!5,colframe=blue!50!black]
\textbf{最後に}：

本ガイドは「UAPは反重力で飛んでいる」と主張するものではありません。
WB理論が予測する飛行特性を正直に列挙し、
公式に報告されたUAPの特性と比較したものです。

一致は示唆的ですが、一致だけでは証拠にはなりません。
科学は実験で進みます。
\$10,000 の実験が、この議論全体の価値を決めます。
\end{tcolorbox}

\end{document}
